\documentclass{article}
\usepackage[utf8]{inputenc}
\usepackage[paperheight=57.15 cm,paperwidth=101.61cm,margin=1in,top=0.7in,heightrounded]{geometry}
\usepackage{xeCJK}
\usepackage{tikz}
\usetikzlibrary{calc}
\usepackage{amsmath}
\usepackage[framemethod=TikZ]{mdframed}
\usepackage{anyfontsize}
\usepackage{sectsty}
\usepackage{graphicx}
\usepackage{fancyhdr}
\usepackage{enumitem}

% Please set fonts by yourself 
\setCJKmainfont[AutoFakeBold=0.3]{cwTeXMing}
\setmainfont{Times New Roman}
\newfontfamily\titlefont{Nunito-Bold.ttf}
\setCJKfamilyfont{kai}{cwTeXKai}


% Poster information
\title{A novel least-squares finite element method for solving convection-dominated problems}
\author{Xxx-Xxx Hsieh (\CJKfamily{kai}{謝XX}),  Department of Mathematics, National Cheng Kung University, Taiwan  (twsiam2020@gmail.com)\\[0.2em] Advisor: Prof. Xxx-Xxx Yang (\CJKfamily{kai}{楊XX})}



\XeTeXlinebreaklocale "zh"
\XeTeXlinebreakskip = 0pt plus 1pt
\input{twsiam.tex}


\begin{document}
\pagestyle{fancy}
\fancyhead{}
\fancyfoot{}
\fancyhead[R]{\textcolor{headgray}{\fontsize{24}{24}\selectfont National Cheng Kung University, Tainan, Taiwan, Jul 24, 2020}}
\fancyhead[L]{\textcolor{headgray}{\fontsize{24}{24}\selectfont 第八屆台灣工業與應用數學會年會}}
\renewcommand{\headrulewidth}{0pt}
\begin{tikzpicture}[overlay,remember picture]
    \draw [line width=2pt,rounded corners=20pt,color=border]
        ($ (current page.north west) + (1cm,-1.5cm) $)
        rectangle
        ($ (current page.south east) + (-1cm,1.5cm) $);
\end{tikzpicture}
\maketitle
\noindent
\begin{center}
\begin{minipage}[t]{0.25\textwidth}
  \begin{myframe}
  \input{contents/01_abstract.tex}
  \end{myframe}
  \begin{myframe}
  \input{contents/04_discussion.tex}
  \end{myframe}
\end{minipage}\hspace{-1em}
\begin{minipage}[t]{0.5\textwidth}
  \begin{myframe}
  \input{contents/02_problem.tex}  
  \end{myframe}
  \begin{myframe}
  \input{contents/05_results.tex}  
  \end{myframe}
\end{minipage}\hspace{-1em}
\begin{minipage}[t]{0.25\textwidth}
  \begin{myframe}
  \input{contents/03_conclusions.tex}  
  \end{myframe}
  \begin{myframe}
  \input{contents/06_references.tex}  
  \end{myframe}
\end{minipage}
\end{center}
\end{document}
